%\documentclass{article}  
\documentclass[12pt,a4paper,openany,fleqn]{book} %声明文档类型

\usepackage{lipsum}
\usepackage[colorlinks]{hyperref}
\renewcommand{\contentsname}{\centerline{目录}}

\usepackage{xeCJK} % 中文
\usepackage{listings} % 插入代码
\usepackage{xcolor} % 代码颜色

%
% 插入代码
%
\lstset{language=C++}%这条命令可以让LaTeX 排版时将C++ 键字突出显示
\lstset{breaklines}%这条命令可以让LaTeX自动将长的代码行换行排版
\lstset{extendedchars=false}% 这一条命令可以解决代码跨页时，章节标题，页眉等汉字不显示的问题
\lstset{
    numbers=left,
    numberstyle= \tiny,
    keywordstyle= \color{ blue!70},
    commentstyle= \color{red!50!green!50!blue!50},
    frame=shadowbox, % 阴影效果
    rulesepcolor= \color{ red!20!green!20!blue!20} ,
    escapeinside=``, % 英文分号中可写入中文
    xleftmargin=2em,xrightmargin=2em, aboveskip=1em,
    framexleftmargin=2em,
}

\setlength\parskip{\baselineskip}% 增加空行%
\setcounter{tocdepth}{4}%控制目录2 级
\setcounter{secnumdepth}{5}%题号显示层级深度

\begin{document}  
\tableofcontents

%
%
%
\chapter{rcu}
\section{rcu}
RCU(Read-Copy Update)，是 Linux 中比较重要的一种同步机制。顾名思义就是“读，拷贝更新”，再直白点是“随意读，但更新数据的时候，需要先复制一份副本，在副本上完成修改，再一次性地替换旧数据”。这是 Linux 内核实现的一种针对“读多写少”的共享数据的同步机制。

不同于其他的同步机制，它允许多个读者同时访问共享数据，而且读者的性能不会受影响（“随意读”），读者与写者之间也不需要同步机制（但需要“复制后再写”），但如果存在多个写者时，在写者把更新后的“副本”覆盖到原数据时，写者与写者之间需要利用其他同步机制保证同步。

RCU 的一个典型的应用场景是链表，在 Linux kernel 中还专门提供了一个头文件（include/linux/rculist.h），提供了利用 RCU 机制对链表进行增删查改操作的接口。

\end{document}  